% Budget: 1.5 pages, ~500 words
\chapter{Introduction}
\label{ch:introduction}

Backlog
% Moved from Positionality on 6 Nov, Celestin's own text, unchanged.

A certain fascination for the topic of cults is not uncommon. Outside of scholarly work, it is a topic that has drawn a lot of attention from mainstream media, especially at the turn of the millennium. The Johnstown Massacre, the Heaven's Gate or Ordre du Temple Solaire suicides, siege were all events that received a lot of media coverage, and they have been the subject of many documentaries and media pieces.

There seems to be a convergence of aspect that turn them into perfect attention-grabbing objects – for better of for worse: they are often violent, tragic, mysterious, and gravitate around beliefs that are often considered extreme or unusual. More interestingly they also touch on fundamental questions about how we relate to each other socially, and how we relate to the world around us.

They also touch on fundamental questions about how we relate one another in a society, and to the world around us: What are a acceptable beliefs? What are acceptable social strucures of powers? Who decides what is acceptable? What is the role of the state in regulating beliefs and social structures?

These questions prove to be relevant when trying to legislate or regulate cults, and they are also relevant when trying to understand how people relate to new technologies, such as AI.

Another aspect that seem to have taken a lot of space is the semantic shift of the word to a an 'adjective' that describes a certain type of social relationship, rather than a noun that describes a specific group. Here we think of 'cult movies', 'cult following', 'cult of personality', and so on. Which can often take a tongue-in-cheek tone, therefore diltuing the negative connotation of the word.

% Headings from v1 (thesis/03_Content/1_Introduction.tex), added 20 Nov; no text
\section{Motivation}
\label{sec:motivation}

\subsection{An unstable word}
\label{subsec:unstable_word}

\subsection{Legitimacy is legally and culturally constructed}
\label{subsec:legitimacy}

\subsection{Cults leave the spiritual paradigm}
\label{subsec:spiritual_paradigm}

\subsection{Toward a structural reading}
\label{subsec:structural_reading}

\section{Research aim}
In this project, the aim is to create an understanding of what we collectively understand to be behind the word cult, both as a category and which groups or entities could be conceived as cults. The blurry nature of the term, the negative connotation, and the social and legal implications of the term make \textit{a priori} attempts to define it seem to never fully grasp the complexity and internal contradictions of the term. The aim is to map the different ways people conceive of cults, draw connections, and build my own understanding of what people really mean when they say cult.

\subsection{Scope}
\label{subsec:scope}

\subsection{Language as mediation tool}

Since this paper mainly focuses on language as a mediation tool, it must be kept in mind that the map we are creating is not a map of the reality of cults but a map of how we conceive it through language. The map is therefore a reflection of the different ways people conceive of cults, and the different ways they use language to describe them, or as described by \textcite{liProxyPresumptionSemantic2026}, one should be wary not to fall into \quote{proxy-presumption}: treating embedding distances as if they directly were social constructs. Its point is that without validation, embeddings mix the concept you want with confounds like topic, style and authorship. In this paper, we'll be treating differences in style as meaningful, and not correcting for them in the methodology. The reason is that this is also a study of language differences, and differences in style are treated as meaningful. They may express similar details, but patching for style would introduce its own biases.

The AI trigger
- A main reason for getting into the topic: the stories of AI psychosis and AI sycophancy.
- Listening to people experiencing these, it sounded very similar to what cult members seem to experience.
- The way people relate to AI is also a social relationship.
- An interesting moment to think about cults: we will look for new words and concepts to understand how we socially and anthropologically relate to technologies that are shifting our work and every aspect of our lives.

\section{Outline of work}
\label{sec:outline}
